\documentclass{beamer}
\usetheme{default}
\beamertemplatenavigationsymbolsempty
\usepackage{rembeamer}

\begin{document}

\begin{frame}
  \begin{itemize}
    \item[] 
      Identifying Type 1 and Type 2 errors in a given situation and
      using them to make decisions about inference methods.
  \end{itemize}
\end{frame}

\begin{frame}
  \frametitle{Decision Errors}
  \begin{table}
    \begin{tabular}{l|l|l}
      Truth & Reject NH & Do Not Reject NH \\[1cm]
      \hline
      NH is true & Type 1 Error & Correct! \\[0.5cm]
      AH is true & Correct! & Type 2 Error
    \end{tabular}
  \end{table}
\end{frame}

\begin{frame}
  \frametitle{Decision Errors - Criminal Trial}
  \begin{table}
    \begin{tabular}{l|l|l}
      Truth & Found Guilty & Found Innocent. \\[1cm]
      \hline
      Defendent is Innocent & Type 1 Error & Correct! \\[0.5cm]
      Defendant is Guilty  & Correct! & Type 2 Error
    \end{tabular}
  \end{table}
\end{frame}

\begin{frame}
  \frametitle{Errors and Risks}
  \begin{itemize}
    \item[] False positive risk: avoid Type 1 errors. 
      \vspace{0.5cm}
    \item[] False negative risk: avoid Type 2 errors. 
  \end{itemize}
\end{frame}

\begin{frame}
  \frametitle{Setting $\alpha$}
  \begin{itemize}
    \item[] 
      Default: $\alpha = 0.05$
      \vspace{0.5cm}
    \item[] 
      Avoid Type 1 error: $\alpha = 0.01$. 
      \vspace{0.5cm}
    \item[] 
      Avoid Type 2 error: $\alpha = 0.10$. 
  \end{itemize}
\end{frame}

\begin{frame}
  \frametitle{Hypotheses}
  \begin{displaymath}
    \begin{array}{rcl}
      H_0 & \mapsto & p_a - p_b = 0 \\[2em]
      H_A & \mapsto & p_a - p_b > 0 \\[2em]
      H_A & \mapsto & p_a - p_b \neq 0 
    \end{array}
  \end{displaymath} 
\end{frame}

\begin{frame}
  \frametitle{Normal Distribution}
  \centering
  \begin{tikzpicture}
    \begin{axis}[
        ytick=\empty,
      ]
      \addplot [
        domain=-3:3,
        samples=200,
      ]
      {exp(-x^2 / (2)) / (sqrt(2*pi))};
    \end{axis}
  \end{tikzpicture}
\end{frame}

\begin{frame}
  \frametitle{Normal Distribution}
  \centering
  \begin{tikzpicture}
    \begin{axis}[
        ytick=\empty,
      ]
      \addplot [
        domain=-3:3,
        samples=200,
      ]
      {exp(-x^2 / (2)) / (sqrt(2*pi))};
    \end{axis}
  \end{tikzpicture}
\end{frame}

\begin{frame}
  \frametitle{Two-Sided Tests}
  \begin{itemize}
    \item[] Two-sided tests make significance twice as difficult. 
      \vspace{0.5cm}
    \item[] Making significance difficult lower the chance of a Type 1
    error. 
      \vspace{0.5cm}
    \item[] Two-sided tests have fewer type one errors.
      \vspace{0.5cm}
    \item[] One-sided tests have fewer type two errors. 
  \end{itemize}
\end{frame}

\begin{frame}
  \frametitle{Other Influences}
  \begin{itemize}
    \item[] 
      Sample size. 
      \vspace{0.5cm}
    \item[] 
      SE $ = \sqrt{ \dfrac{\hat{p}_1 (1 - \hat{p}_1)}{n_1} +
      \dfrac{\hat{p}_2 (1 - \hat{p}_2)}{n_2}}$ 
      \vspace{0.5cm}
    \item[] 
      The \emph{power} of a statistical test. 
  \end{itemize}
\end{frame}

\end{document}

